\section{Metodologija}
\label{sec:metodologija}

Istraživanje je eksperimentalno-kvantitativnog tipa. Provodi se empirijska
evaluacija četiri konfiguracije vektorskih indeksa nad istim skupom podataka,
uz precizno mjerenje latencije upita, tačnosti pretrage i memorijskog
opterećenja po konfiguraciji. Poglavlje opisuje eksperimentalno okruženje
i skup podataka, postupak generisanja vektorskih reprezentacija, izgradnju
indeksnih struktura te protokol mjerenja i statističke obrade rezultata.

\subsection{Eksperimentalno okruženje i skup podataka}
\label{sec:okruzenje_podaci}

Svi eksperimenti provode se na prijenosnom računaru Lenovo Legion 5 16IRX9
opremljenom procesorom Intel Core i7-14650HX s 16 jezgara i maksimalnom
taktnom frekvencijom od $5{,}2$\,GHz, $32$\,GB radne memorije i NVMe SSD
pogonom kapaciteta $953{,}87$\,GB. Za akceleraciju enkodiranja dostupna je
grafička kartica NVIDIA GeForce RTX 4070 Laptop GPU. Operativni sistem je
Ubuntu 24.04.4 LTS s jezgrom verzije 6.17.0.

Eksperiment se izvodi u programskom jeziku Python. Koriste se sljedeće
biblioteke i servisi: \foreignlanguage{english}{Qdrant} kao vektorska baza
podataka pokrenuta unutar \foreignlanguage{english}{Docker} kontejnera,
\foreignlanguage{english}{FAISS} kao biblioteka za aproksimativnu pretragu
vektora, \foreignlanguage{english}{sentence-transformers} za generisanje
vektorskih reprezentacija, \foreignlanguage{english}{datasets} za učitavanje
skupa podataka, te \foreignlanguage{english}{NumPy}, \foreignlanguage{english}{SciPy},
\foreignlanguage{english}{pandas} i \foreignlanguage{english}{Matplotlib}
za numeričku obradu, statističku analizu i vizualizaciju rezultata.

Kao skup podataka koristi se \foreignlanguage{english}{MS MARCO v2.1}
(\foreignlanguage{english}{Microsoft MAchine Reading COmprehension}),
standardni benchmark za zadatke pretraživanja informacija koji sadrži upite
u prirodnom jeziku i odgovarajuće pasuse iz stvarnih web
dokumenata~\cite{msmarco2016}. Svaki primjer u skupu podataka sadrži polje
\emph{query} s upitom u prirodnom jeziku te polje \emph{passages} koje
obuhvata skup pasusa, od kojih je svaki opisan tekstualnim sadržajem
(\emph{passage\_text}) i oznakom relevantnosti (\emph{is\_selected}).

Iz train dijela skupa podataka, koji broji oko $809{.}000$ primjera,
ekstrahuju se svi pasusi koji čine korpus za indeksiranje. Iz validation
dijela, koji sadrži oko $101{.}000$ primjera, nasumično se odabira $1{.}000$
upita koji se koriste za sva mjerenja. Odabir se vrši uz fiksiran generator
slučajnih brojeva radi reproducibilnosti eksperimenta. Isti skup od $1{.}000$
upita primjenjuje se na sve četiri konfiguracije indeksa, što je ključan
uvjet za statističku valjanost poređenja.

\subsection{Generisanje vektorskih reprezentacija}
\label{sec:embeddings}

Pasusi i upiti prevode se u numeričke vektore korištenjem modela
\foreignlanguage{english}{\emph{all-mpnet-base-v2}},
koji generira $768$-dimenzionalne vektore optimizirane za mjerenje semantičke
sličnosti između rečenica. Model se primjenjuje odvojeno na korpus pasusa
i skup upita.

Enkodiranje se provodi u serijama od $128$ tekstova istovremeno. Veličina
serije određuje koliko se tekstova obrađuje paralelno, što direktno utiče
na brzinu enkodiranja i iskoristivost memorije grafičke kartice, a vrijednost
$128$ predstavlja kompromis između ta dva faktora koji je primjeren raspoloživom
hardveru.

Dobijeni vektori normalizuju se L2 normalizacijom, kojom se svaki vektor
dijeli sa svojom euklidskom normom tako da rezultujući vektor leži na
jediničnoj sferi. Ovim postupkom euklidska udaljenost i kosinusna sličnost
postaju ekvivalentne mjere rangiranja~\cite{manning2008introduction}, što
omogućava konzistentnu upotrebu \foreignlanguage{english}{FAISS} indeksa
zasnovanog na euklidskoj udaljenosti i \foreignlanguage{english}{Qdrant}
servisa koji koristi kosinusnu sličnost. Ispravnost normalizacije verifikuje
se numeričkom provjerom normi vektora. Generirani vektori pohranjuju se na
disk u binarnom formatu radi ponovnog korištenja u svim eksperimentima.

\subsection{Izgradnja indeksnih struktura}
\label{sec:indeksi}

Testiraju se četiri konfiguracije indeksa podijeljene između dva alata.
\foreignlanguage{english}{Flat}, \foreignlanguage{english}{IVF} i
\foreignlanguage{english}{IVF-PQ} indeksi grade se unutar
\foreignlanguage{english}{FAISS} biblioteke, koja
radi direktno u memoriji procesa putem C++ vezanja s minimalnim sistemskim
troškom. \foreignlanguage{english}{HNSW} indeks gradi se unutar
\foreignlanguage{english}{Qdrant} servisa koji radi kao produkcijski
mikroservis s mrežnim interfejsom (\foreignlanguage{english}{HTTP} i
\foreignlanguage{english}{gRPC}). Ovakva podjela osigurava da svaki indeks
bude testiran u tehnološkom okruženju za koje je optimiziran, što rezultatima
daje praktičnu relevantnost.

Budući da \foreignlanguage{english}{Qdrant} kao mikroservis unosi dodatni
trošak mrežne komunikacije i serijalizacije podataka koji nije prisutan u
\foreignlanguage{english}{FAISS} okruženju, pri svakom eksperimentu mjeri
se i prosječna latencija praznog upita prema \foreignlanguage{english}{Qdrant}
servisu, odnosno latencija samog mrežnog povratnog puta bez stvarne pretrage.
Ta vrijednost bilježi se i oduzima od izmjerenih \foreignlanguage{english}{HNSW}
latencija pri apsolutnom poređenju s \foreignlanguage{english}{FAISS}
konfiguracijama, čime se osigurava da se uspoređuju algoritamske performanse,
a ne arhitekturalni overhead.

\emph{\foreignlanguage{english}{Flat} indeks} provodi iscrpnu pretragu
poređenjem upita sa svakim vektorom u korpusu i garantuje potpunu tačnost.
Služi kao gornja granica za tačnost i donja granica za brzinu te kao referentni
rezultat (\foreignlanguage{english}{\emph{ground truth}}) za računanje tačnosti
preostalih konfiguracija.

\emph{\foreignlanguage{english}{IVF} indeks}
(\foreignlanguage{english}{Inverted File Index}) dijeli vektorski prostor
na $n_{\mathrm{list}}$ klastera algoritmom k-means. Pretraga se ograničava
na $n_{\mathrm{probe}}$ najbližih klastera, čime se smanjuje broj poređenja.
Viša vrijednost $n_{\mathrm{probe}}$ povećava tačnost ali i latenciju, što
daje kontinuirani kompromis između ta dva faktora. Izgradnja indeksa zahtijeva
fazu treniranja k-means algoritmom nad podskupom korpusa.

\emph{\foreignlanguage{english}{IVF-PQ} indeks} kombinuje
\foreignlanguage{english}{IVF} klasterizaciju s
\foreignlanguage{english}{Product Quantization} kompresijom vektora opisanom
u sekciji~\ref{sec:vektorska_kvantizacija}. Kompresija drastično smanjuje
memorijski otisak indeksa na račun aproksimativnosti pri računanju udaljenosti.

\emph{\foreignlanguage{english}{HNSW} indeks} gradi
višeslojnu grafovsku strukturu opisanu u
sekciji~\ref{sec:vektorsko_indeksiranje}.
Parametar $M$ određuje maksimalni broj bidirekcionalnih veza koje svaki
čvor održava u grafu, dok parametar $ef\_\mathrm{construct}$ određuje
koliko kandidata za susjede algoritam razmatra pri izgradnji svake veze, viša
vrijednost $ef\_\mathrm{construct}$ daje precizniji graf uz sporiji
proces izgradnje indeksa. Pri pretrazi, parametar $ef\_\mathrm{search}$
kontroliše koliko kandidata se prati kroz graf, čime se reguliše kompromis
između tačnosti i latencije upita.

\subsection{Mjerenje, metrike i statistička analiza}
\label{sec:mjerenje_statistika}

Za svaku od četiri konfiguracije indeksa provodi se $1{.}000$ upita,
pri čemu se svaki upit ponavlja $5$ puta, što daje ukupno $5{.}000$
mjerenja latencije po konfiguraciji. Latencija se mjeri za svaki upit
pojedinačno korištenjem preciznog sistemskog mjerača vremena, a rezultati
se pohranjuju za kasniju obradu. Rezultati prvog ponavljanja koriste se
za računanje tačnosti pretrage, a svih pet ponavljanja za analizu latencije.

Svaka konfiguracija evaluira se prema trima mjernim metrikama.
\emph{Latencija upita} mjeri se u milisekundama za svaki upit. Iz $5{.}000$
mjerenja po konfiguraciji izračunava se srednja vrijednost i $95.$ percentil
($P_{95}$), koji predstavlja latenciju u najlošijih $5\%$ slučajeva i
relevantan je za procjenu ponašanja sistema u uvjetima opterećenja.
\emph{\foreignlanguage{english}{Recall@10}} mjeri koji udio od $10$ rezultata
koje je aproksimativni indeks vratio odgovara $10$ rezultata
\foreignlanguage{english}{Flat} indeksa uzetog kao referenca. Vrijednost
bliska $1{,}0$ znači da aproksimativna metoda pronalazi gotovo iste rezultate
kao i iscrpna pretraga. \emph{Memorijski otisak} mjeri se kao veličina
serijaliziranog indeksa na disku u megabajtima.

Za svaku od četiri konfiguracije računaju se deskriptivni statistički
pokazatelji latencija: srednja vrijednost, medijana, standardna devijacija,
minimum, maksimum te $P_{95}$ i $P_{99}$ percentili. Za provjeru statističke
značajnosti razlika između latencija dviju konfiguracija primjenjuje se
Wilcoxon \foreignlanguage{english}{\emph{signed-rank}} test~\cite{demsar2006},
neparametrijski test koji ne pretpostavlja normalnu distribuciju podataka
i prikladan je za distribucije latencija koje su tipično asimetrične s dugim
desnim repom. Svaka od tri hipoteze testira se na razini značajnosti
$p < 0{,}05$. Pored statističkih testova, za svaki indeks konstruišu se
krivulje kompromisa između tačnosti i latencije variranjem parametra pretrage
($n_{\mathrm{probe}}$ za \foreignlanguage{english}{IVF} i
\foreignlanguage{english}{IVF-PQ}, $ef\_\mathrm{search}$ za
\foreignlanguage{english}{HNSW}), čime se vizualizira raspon mogućih
kompromisa dostupnih za svaki indeks.
